\bta{2022}






\begin{enumerate}
	%\renewcommand{\labelenumi}{\arabic{enumi}.}
	% A(\Alph) a(\alph) I(\Roman) i(\roman) 1(\arabic)
	%设定全局标号series=example	%引用全局变量resume=example
	%[topsep=-0.3em,parsep=-0.3em,itemsep=-0.3em,partopsep=-0.3em]
	%可使用leftmargin调整列表环境左边的空白长度 [leftmargin=0em]
	\item
	\begin{enumerate}
		%\renewcommand{\labelenumi}{\arabic{enumi}.}
		% A(\Alph) a(\alph) I(\Roman) i(\roman) 1(\arabic)
		%设定全局标号series=example	%引用全局变量resume=example
		%[topsep=-0.3em,parsep=-0.3em,itemsep=-0.3em,partopsep=-0.3em]
		%可使用leftmargin调整列表环境左边的空白长度 [leftmargin=0em]
		\item
		$f(x)=\arcsin \left(\frac{4}{\pi} x-1\right)$, 求 $f[f(x)]$ 的定义域 \xzanswer{A} 
		
		
		
	\fourchoices
	{$\left[\frac{\pi}{4}, \frac{\pi}{2}\right]$}
	{$\left[\frac{\pi}{4}, \frac{\pi}{2}\right]$}
	{$\left[\frac{\pi}{8}, \frac{\pi}{4}\right]$}
	{$\left[0, \frac{\pi}{2}\right]$}
	

	\item 
	求极限 $\lim\limits _{x \rightarrow 0} \tan ^{2} x\left(\frac{1}{\sin x}+\frac{1}{x^{2}}\right)=$ \xzanswer{B} 
	
	
	
\fourchoices
{$ 0 $}
{$ 1 $}
{$ 2 $}
{不存在}

\item 
若函数 $g(x)$ 二阶可导, 且满足 $g(0)=g^{\prime}(0)=0, f(x)=\left\{\begin{array}{cl}\frac{g(x)}{x} & x \neq 0 \\ 0 & x=0\end{array}\right.$, 讨论$f(x)$ 和 $f^{\prime}(x)$ 在 $x=0$ 处的连续性 \xzanswer{A} 

	
\fourchoices
{$f(x)$ 连续, $f^{\prime}(x)$ 连续}
{$f(x)$ 连续, $f^{\prime}(x)$ 不连续, $f^{\prime}(0)$ 存在}
{$f(x)$ 连续, $f^{\prime}(0)$ 不存在}
{$f(x)$ 不连续	}


\item 
$f(x)=\lim\limits _{n \rightarrow+\infty} \frac{e^{x}+x^{2 n}}{2-x^{2 n}}$, 则 $f(x)$ 的间断点为 \xzanswer{C} 




\fourchoices
{$-1$}
{$ 1 $}
{$ \pm 1 $}
{无间断点}


\item 
直线: $\frac{x-1}{2}=\frac{y-2}{-2}=\frac{z-3}{-1}$ 和直线 $\left\{\begin{array}{l}x+2 y-3=0 \\ 3 y+z+2=0\end{array}\right.$ 间夹角的余弦值 $\cos \theta$为 \xzanswer{D} 



\fourchoices
{$\frac{1}{\sqrt{8}}$}
{$\frac{1}{\sqrt{10}}$}
{$\frac{1}{\sqrt{12}}$}
{$\frac{1}{\sqrt{14}}$}

\item 
正项级数 $\sum\limits_{n=1}^{+\infty} a_{n}$ 收敛, 则下列级数：
①$\sum\limits_{n=1}^{+\infty}(-1)^{n} a_{n}$，
②$\sum\limits_{n=1}^{+\infty}\left(a_{n}\right)^{2}$，
③$\sum\limits_{n=1}^{+\infty} \sin \left(a_{n}\right)$，
④$\sum\limits_{n=1}^{+\infty} \ln \left(1+a_{n}\right)$
中一定收敛的级数个数为 \xzanswer{D} 



\fourchoices
{$ 1 $}
{$ 2 $}
{$ 3 $}
{$ 4 $}


\item 
判断曲线: $y=x \sin \frac{1}{x}$ 的渐近线 \xzanswer{B} 

\fourchoices
{只有垂直渐近线}
{只有水平渐近线}
{有垂直渐近线和水平渐近线}
{只有斜渐近线}


\item 
函数 $f(x, y)=\frac{x^{2} y^{2}}{x^{2} y^{2}+(x-y)^{2}}$, 则 $f(x, y)$ 在点 $(0,0)$ 处 \xzanswer{B} 

\fourchoices
{累次极限都存在, 重极限存在}
{累次极限都存在, 重极限不存在}
{累次极限存在其中一个, 重极限不存在}
{至少一个累次极限不存在}

\item 
函数 $z(x, y)$ 满足 $\frac{x}{z}=\ln \frac{z}{y}$, 则 $d z=$ \xzanswer{B} 

\fourchoices
{$\frac{z(z d x+y d y)}{z(x+y)}$}
{$\frac{z(y d x+z d y)}{y(x+z)}$}
{$ \frac{z(z d x+x d y)}{y(x+z)}$}
{$\frac{z(x d x+y d y)}{z(x+y)}$}

\item 
下列广义积分收敛的是 \xzanswer{B} 


\fourchoices
{$\int_{0}^{+\infty} \frac{1}{x\left(e^{\frac{x}{2}}-1\right)} d x$}
{$\int_{0}^{1}(\ln x)^{100} d x$}
{$ \int_{2}^{+\infty} \frac{1}{x \ln x} d x$}
{$\int_{0}^{+\infty} \frac{1}{1+x+\sin x} d x$}







\end{enumerate}


\banswer{
\begin{enumerate}
	%\renewcommand{\labelenumi}{\arabic{enumi}.}
	\item
	\textbf{A}。先求 $f$ 的定义域：由 $-1\leqslant\frac{4}{\pi}x-1\leqslant1$ 得 $0\leqslant x\leqslant\frac{\pi}{2}$。
	要使 $f[f(x)]$ 有意义，除 $x\in\left[0,\frac{\pi}{2}\right]$ 外还需 $f(x)\in\left[0,\frac{\pi}{2}\right]$。
	而 $f(x)=\arcsin\left(\frac{4}{\pi}x-1\right)\leqslant\frac{\pi}{2}$ 自动成立，
	$f(x)\geqslant0\iff\frac{4}{\pi}x-1\geqslant0\iff x\geqslant\frac{\pi}{4}$。
	故定义域为 $\left[\frac{\pi}{4},\frac{\pi}{2}\right]$。

	\item
	\textbf{B}。拆成两项。由 $\tan x\sim x$，
	\[
	\frac{\tan^2x}{\sin x}\sim\frac{x^2}{x}=x\to0,
	\qquad
	\frac{\tan^2x}{x^2}\to1 ,
	\]
	故原极限 $=0+1=1$。

	\item
	\textbf{A}。由 $g(0)=g'(0)=0$ 及二阶可导，Taylor 展开
	$g(x)=\frac{g''(0)}{2}x^2+o\left(x^2\right)$，$g'(x)=g''(0)x+o(x)$。

	$f$ 连续：$\lim\limits_{x\to0}\frac{g(x)}{x}=g'(0)=0=f(0)$。

	$f'(0)$ 存在：$f'(0)=\lim\limits_{x\to0}\frac{f(x)-f(0)}{x}=\lim\limits_{x\to0}\frac{g(x)}{x^2}=\frac{g''(0)}{2}$。

	$f'$ 连续：$x\neq0$ 时 $f'(x)=\frac{xg'(x)-g(x)}{x^2}$，而
	\[
	xg'(x)-g(x)=g''(0)x^2-\frac{g''(0)}{2}x^2+o\left(x^2\right)=\frac{g''(0)}{2}x^2+o\left(x^2\right),
	\]
	故 $\lim\limits_{x\to0}f'(x)=\frac{g''(0)}{2}=f'(0)$。故 $f$ 与 $f'$ 都连续。

	\item
	\textbf{C}。按 $|x|$ 与 $1$ 的大小分段取极限：
	\[
	f(x)=
	\begin{cases}
	\dfrac{e^x}{2}, & |x|<1\quad\left(x^{2n}\to0\right),\\[6pt]
	e^x+1, & |x|=1\quad\left(x^{2n}\equiv1\right),\\[6pt]
	-1, & |x|>1\quad\left(x^{2n}\to+\infty\right).
	\end{cases}
	\]
	在 $x=1$ 处左极限 $\frac{e}{2}$、右极限 $-1$、函数值 $e+1$，三者互不相等；
	$x=-1$ 处同理。故间断点为 $x=\pm1$。

	\item
	\textbf{D}。第一条直线方向 $\bm{s}_1=(2,-2,-1)$，$\left|\bm{s}_1\right|=3$。
	第二条直线是两平面的交线，方向
	$\bm{s}_2=(1,2,0)\times(0,3,1)=(2,-1,3)$，$\left|\bm{s}_2\right|=\sqrt{14}$。故
	\[
	\cos\theta=\frac{\left|\bm{s}_1\cdot\bm{s}_2\right|}{\left|\bm{s}_1\right|\left|\bm{s}_2\right|}
	=\frac{|4+2-3|}{3\sqrt{14}}=\frac{3}{3\sqrt{14}}=\frac{1}{\sqrt{14}} .
	\]

	\item
	\textbf{D}。由 $\sum a_n$ 收敛且 $a_n>0$ 知 $a_n\to0$，从而存在 $N$ 使 $n>N$ 时 $0<a_n<1$。
	① $\left|(-1)^na_n\right|=a_n$，绝对收敛；
	② $n>N$ 时 $a_n^2<a_n$；
	③ $0<\sin a_n\leqslant a_n$；
	④ $0<\ln\left(1+a_n\right)\leqslant a_n$。
	四者都可由比较判别法得到收敛，故个数为 $4$。

	\item
	\textbf{B}。$x\to\infty$ 时 $x\sin\frac1x=\frac{\sin(1/x)}{1/x}\to1$，得水平渐近线 $y=1$；
	$x\to0$ 时 $\left|x\sin\frac1x\right|\leqslant|x|\to0$ 为有界量，不趋于 $\infty$，故无垂直渐近线。
	又 $\frac{y}{x}=\sin\frac1x\to0$，斜渐近线退化为上述水平渐近线。故只有水平渐近线。

	\item
	\textbf{B}。两个累次极限都存在且为 $0$：
	$\lim\limits_{y\to0}f(x,y)=\frac{0}{x^2}=0\ (x\neq0)$，再令 $x\to0$ 得 $0$；由对称性另一个也是 $0$。
	但重极限不存在：沿 $y=x$ 有 $f=\frac{x^4}{x^4+0}=1\to1$，
	沿 $y=0$ 有 $f\equiv0\to0$，两者不等。

	\item
	\textbf{B}。令 $F(x,y,z)=\frac{x}{z}-\ln z+\ln y$，则
	$F_x=\frac1z$，$F_y=\frac1y$，$F_z=-\frac{x}{z^2}-\frac1z=-\frac{x+z}{z^2}$。故
	\[
	dz=-\frac{F_xdx+F_ydy}{F_z}
	=\frac{\frac1zdx+\frac1ydy}{\frac{x+z}{z^2}}
	=\frac{z\,dx+\frac{z^2}{y}dy}{x+z}
	=\frac{z(y\,dx+z\,dy)}{y(x+z)} .
	\]

	\item
	\textbf{B}。A 中 $x\to0^{+}$ 时 $e^{x/2}-1\sim\frac{x}{2}$，被积函数 $\sim\frac{2}{x^2}$，发散；
	C 中 $\int_2^{+\infty}\frac{dx}{x\ln x}=\left[\ln\ln x\right]_2^{+\infty}=+\infty$，发散；
	D 中 $x\to+\infty$ 时被积函数 $\sim\frac1x$，发散。
	B 中令 $x=e^{-t}$，
	\[
	\int_0^1\ln^{100}x\,dx=\int_0^{+\infty}t^{100}e^{-t}dt=\Gamma(101)=100! ,
	\]
	收敛。

	\par\vspace{1mm}\noindent
	\textbf{注：}{\kaishu
	第 (1) 题现有题面中选项 A 与 B 完全相同，应是誊录时的笔误，正确答案为
	$\left[\frac{\pi}{4},\frac{\pi}{2}\right]$。
	}\par
\end{enumerate}
}



\item 
$u(x)=\left\{\begin{array}{ll}1,& x>0 \\ 0,& x \leq 0\end{array}\right.$, 求极限 $\lim\limits _{x \rightarrow 0}\left[\frac{1+e^{\frac{1}{x}}}{1-e^{\frac{2}{x}}}+u(x)\right]$。

\banswer{
	$x=0$ 是 $e^{1/x}$ 的本性间断点，需分别求左、右极限。

	$x\to0^{-}$ 时 $\frac1x\to-\infty$，故 $e^{1/x}\to0$，$e^{2/x}\to0$，且 $u(x)=0$：
	\[
	\lim_{x\to0^{-}}\left[\frac{1+e^{1/x}}{1-e^{2/x}}+u(x)\right]=\frac{1+0}{1-0}+0=1 .
	\]

	$x\to0^{+}$ 时 $\frac1x\to+\infty$，分子分母同除以 $e^{2/x}$：
	\[
	\frac{1+e^{1/x}}{1-e^{2/x}}=\frac{e^{-2/x}+e^{-1/x}}{e^{-2/x}-1}\longrightarrow\frac{0+0}{0-1}=0 ,
	\]
	而 $u(x)=1$，故右极限 $=0+1=1$。

	左右极限相等，所以
	\[
	\lim_{x\to0}\left[\frac{1+e^{1/x}}{1-e^{2/x}}+u(x)\right]=1 .
	\]
}


\item 
函数 $y(x)$ 满足微分方程 $y^{\prime \prime}=e^{2 y}+e^{y}$, 且 $y(0)=0, y^{\prime}(0)=2$, 求 $y(x)$。


\banswer{
	方程不显含 $x$，作降阶代换 $p=y^{\prime}$，视 $p$ 为 $y$ 的函数，则
	$y^{\prime\prime}=\frac{dp}{dx}=p\frac{dp}{dy}$，原方程化为
	\[
	p\,dp=\left(e^{2y}+e^{y}\right)dy
	\quad\Longrightarrow\quad
	\frac{p^2}{2}=\frac{e^{2y}}{2}+e^{y}+C .
	\]
	代入 $y(0)=0$、$p(0)=y^{\prime}(0)=2$ 得 $2=\frac12+1+C$，故 $C=\frac12$，于是
	\[
	p^2=e^{2y}+2e^{y}+1=\left(e^{y}+1\right)^2 .
	\]
	因 $p(0)=2=e^{0}+1>0$，由连续性取正号：$\frac{dy}{dx}=e^{y}+1$。分离变量，
	\[
	\int\frac{dy}{e^{y}+1}=\int\frac{e^{-y}\,dy}{1+e^{-y}}=-\ln\left(1+e^{-y}\right)=x+C_1 .
	\]
	由 $x=0$、$y=0$ 得 $C_1=-\ln2$，故 $\ln\left(1+e^{-y}\right)=\ln2-x$，即
	$1+e^{-y}=2e^{-x}$，$e^{-y}=\frac{2-e^{x}}{e^{x}}$。所以
	\[
	y=\ln \frac{e^{x}}{2-e^{x}}\qquad\left(x<\ln2\right).
	\]
	（验算：$y^{\prime}=\frac{2}{2-e^{x}}$，$y^{\prime\prime}=\frac{2e^{x}}{\left(2-e^{x}\right)^2}$，
	而 $e^{2y}+e^{y}=\frac{e^{2x}+e^{x}\left(2-e^{x}\right)}{\left(2-e^{x}\right)^2}
	=\frac{2e^{x}}{\left(2-e^{x}\right)^2}$，相等。）
}


\item 
设函数 $y(x)$ 满足参数方程 $\left\{\begin{array}{l}x=1+t^{2} \\ y=\cos t\end{array}\right.$，求下列极限：
\begin{enumerate}
	%\renewcommand{\labelenumi}{\arabic{enumi}.}
	% A(\Alph) a(\alph) I(\Roman) i(\roman) 1(\arabic)
	%设定全局标号series=example	%引用全局变量resume=example
	%[topsep=-0.3em,parsep=-0.3em,itemsep=-0.3em,partopsep=-0.3em]
	%可使用leftmargin调整列表环境左边的空白长度 [leftmargin=0em]
	\item

 $ \frac{d y}{d x}$和$ \frac{d^{2} y}{d x^{2}}$；
\item 
$\lim\limits _{x \rightarrow 1^{+}} \frac{d y}{d x} $和$ \lim\limits _{x \rightarrow 1^{+}} \frac{d^{2} y}{d x^{2}}$。
\end{enumerate}

\banswer{
	\begin{enumerate}
		%\renewcommand{\labelenumi}{\arabic{enumi}.}
		% A(\Alph) a(\alph) I(\Roman) i(\roman) 1(\arabic)
		%设定全局标号series=example	%引用全局变量resume=example
		%[topsep=-0.3em,parsep=-0.3em,itemsep=-0.3em,partopsep=-0.3em]
		%可使用leftmargin调整列表环境左边的空白长度 [leftmargin=0em]
		\item
		由参数方程求导公式，$\frac{dx}{dt}=2t$，$\frac{dy}{dt}=-\sin t$，故
		\[
		\frac{dy}{dx}=\frac{dy/dt}{dx/dt}=-\frac{\sin t}{2 t}\qquad(t\neq0).
		\]
		再对 $x$ 求导（先对 $t$ 求导，再除以 $\frac{dx}{dt}$）：
		\[
		\frac{d^{2} y}{d x^{2}}=\frac{\dfrac{d}{dt}\left(-\dfrac{\sin t}{2t}\right)}{\dfrac{dx}{dt}}
		=\frac{-\dfrac{t\cos t-\sin t}{2t^{2}}}{2t}
		=-\frac{t \cos t-\sin t}{4 t^{3}} .
		\]
		
		\item 
		$x\to1^{+}$ 即 $t\to0$（$t\neq0$）。由 $\sin t\sim t$，
		\[
		\lim_{x\to1^{+}}\frac{dy}{dx}=\lim_{t\to0}\left(-\frac{\sin t}{2t}\right)=-\frac{1}{2} .
		\]
		对第二式用 Taylor 展开：
		$t\cos t=t-\frac{t^{3}}{2}+o\left(t^{3}\right)$，$\sin t=t-\frac{t^{3}}{6}+o\left(t^{3}\right)$，故
		$t\cos t-\sin t=-\frac{t^{3}}{3}+o\left(t^{3}\right)$，于是
		\[
		\lim_{x\to1^{+}}\frac{d^{2} y}{d x^{2}}
		=\lim_{t\to0}\left(-\frac{-\frac{t^{3}}{3}+o\left(t^{3}\right)}{4t^{3}}\right)
		=\frac{1}{12} .
		\]
		
	\end{enumerate}
	
	
}



\item 
证明: $u(x, y)=f(x) g(y)$ 成立的充分必要条件是 $u \frac{\partial^{2} u}{\partial x \partial y}=\frac{\partial u}{\partial x} \frac{\partial u}{\partial y}$。
\banswer{
	\textbf{必要性：} 若 $u(x,y)=f(x)g(y)$，则
	$\frac{\partial u}{\partial x}=f'(x)g(y)$，$\frac{\partial u}{\partial y}=f(x)g'(y)$，
	$\frac{\partial^{2} u}{\partial x\partial y}=f'(x)g'(y)$。
	于是 $u\frac{\partial^{2} u}{\partial x\partial y}=f(x)g(y)\cdot f'(x)g'(y)=\frac{\partial u}{\partial x}\frac{\partial u}{\partial y}$。\\[6pt]
	\textbf{充分性（原命题不成立）：} 即使
	$u\frac{\partial^{2}u}{\partial x\partial y}=\frac{\partial u}{\partial x}\frac{\partial u}{\partial y}$
	处处成立，全局也未必有 $u(x,y)=f(x)g(y)$。关键在于乘以 $u$ 的因子化可以抵消零点的干扰。

	\textbf{反例}：取光滑截断函数 $\varphi_-,\varphi_+\in C_c^\infty(\mathbb{R})$，
	$\operatorname{supp}\varphi_-\subset(-2,-1)$，$\operatorname{supp}\varphi_+\subset(1,2)$，
	且均不恒为零。令
	\[
	u(x,y)=\varphi_-(x)+y\,\varphi_+(x),\qquad (x,y)\in\mathbb{R}^2.
	\]
	对任意 $(x,y)$，$u_x=\varphi_-'(x)+y\varphi_+'(x)$，
	$u_y=\varphi_+(x)$，$u_{xy}=\varphi_+'(x)$。
	计算 $u\,u_{xy}-u_xu_y=[\varphi_-(x)+y\varphi_+(x)]\varphi_+'(x)-[\varphi_-'(x)+y\varphi_+'(x)]\varphi_+(x)
	=\varphi_-(x)\varphi_+'(x)-\varphi_-'(x)\varphi_+(x)$。
	由于 $\varphi_-$ 与 $\varphi_+$ 的支集不相交，对任意 $x$ 至少有一个为零，从而差恒为零，
	故 PDE 对任意 $(x,y)$ 成立。然而 $u(x,y)$ 不能写成 $f(x)g(y)$：
	若可分解，则固定 $x_1\in(-2,-1)$ 使 $\varphi_-(x_1)\neq0$，固定 $x_2\in(1,2)$ 使 $\varphi_+(x_2)\neq0$，
	考虑 $u(x_1,y)=\varphi_-(x_1)$（与 $y$ 无关），而 $u(x_2,y)=y\varphi_+(x_2)$（线性依赖 $y$）；
	同时要求 $u(x,y)=f(x)g(y)$，则从 $u(x_1,y)$ 得 $g(y)$ 在支撑上为常数，
	从 $u(x_2,y)$ 却要求 $g(y)$ 非平凡依赖 $y$，矛盾。
	故原命题的充分性在全局意义下为假，其成立需要额外的非零性假设。

	\textbf{修正定理（非零情形）：} 设 $u$ 在矩形区域
	$I\times J\subset\mathbb{R}^2$ 上连续可微，且 $u(x,y)\neq0$ 于 $I\times J$。
	则 $u\,u_{xy}=u_xu_y$ 成立的充要条件是 $u(x,y)=f(x)g(y)$，其中 $f,g$ 为非零可微函数。

	\textbf{修正定理的证明：}
	由 $u\neq0$ 可作除法。考虑
	\[
	\frac{\partial}{\partial y}\!\left(\frac{u_x}{u}\right)
	=\frac{u\,u_{xy}-u_xu_y}{u^{2}}=0,
	\]
	故 $\dfrac{u_x}{u}$ 与 $y$ 无关。于是在 $I\times J$ 上
	\[
	\frac{u_x}{u}=A(x),\qquad \forall\,(x,y)\in I\times J.
	\]
	对固定的 $y$ 关于 $x$ 积分得
	$\ln|u(x,y)|=\int_{x_0}^{x} A(t)\,dt + h(y)$，
	其中 $h(y)$ 为仅依赖于 $y$ 的积分常数。从而
	\[
	u(x,y)=\pm e^{\int_{x_0}^{x}A(t)\,dt}\cdot e^{h(y)}\equiv f(x)\,g(y),
	\]
	其中 $f(x)=\pm e^{\int_{x_0}^{x}A(t)\,dt}$，$g(y)=e^{h(y)}$，均非零。
}




\item 
计算半圆 $(x-1)^{2}+y^{2}=1$，$ y>0$, 绕 $y$ 轴旋转所得到的旋转体的体积。


\banswer{
	旋转体由上半圆盘 $D:\ (x-1)^2+y^2\leqslant1,\ y\geqslant0$ 绕 $y$ 轴旋转而成，
	$D$ 在 $x$ 轴上的跨度是 $0\leqslant x\leqslant2$，故用\textbf{柱壳法}：
	\[
	V=\int_0^2 2\pi x\cdot y_{\text{上}}(x)\,dx=2\pi\int_0^2 x\sqrt{1-(x-1)^2}\,dx .
	\]
	令 $u=x-1$，
	\[
	V=2\pi\int_{-1}^{1}(u+1)\sqrt{1-u^2}\,du
	=2\pi\underbrace{\int_{-1}^{1}u\sqrt{1-u^2}\,du}_{=\,0\ (\text{奇函数})}
	+2\pi\int_{-1}^{1}\sqrt{1-u^2}\,du=2\pi\cdot\frac{\pi}{2}=\pi^{2} .
	\]
	（其中 $\int_{-1}^{1}\sqrt{1-u^2}\,du$ 是半径为 $1$ 的半圆面积 $\frac{\pi}{2}$。）

	\par\vspace{1mm}\noindent
	\textbf{注：}{\kaishu
	也可直接用 Pappus（古尔丁）定理：半圆面积 $A=\frac{\pi}{2}$，
	其形心横坐标由对称性为 $\bar{x}=1$，故 $V=2\pi\bar{x}A=2\pi\cdot1\cdot\frac{\pi}{2}=\pi^{2}$。
	}\par
}


\item 
曲线 $C: \frac{x^{2}}{4}+\frac{y^{2}}{9}=1$, 计算积分
\[ I=\oint_{C} e^{x y}\{[y \sin (x y)+\cos (x+y)] d x+[x \sin (x y)+\cos (x+y)] d y\} \]

\banswer{
	记 $P=e^{xy}\left[y\sin(xy)+\cos(x+y)\right]$，$Q=e^{xy}\left[x\sin(xy)+\cos(x+y)\right]$，
	$D$ 为 $C$ 所围的椭圆域。分别求偏导：
	\begin{align*}
	\frac{\partial P}{\partial y}&=e^{xy}\Bigl[xy\sin(xy)+\sin(xy)+xy\cos(xy)+x\cos(x+y)-\sin(x+y)\Bigr],\\
	\frac{\partial Q}{\partial x}&=e^{xy}\Bigl[xy\sin(xy)+\sin(xy)+xy\cos(xy)+y\cos(x+y)-\sin(x+y)\Bigr].
	\end{align*}
	两式只差最后一组，故
	\[
	\frac{\partial Q}{\partial x}-\frac{\partial P}{\partial y}=e^{xy}(y-x)\cos(x+y)\not\equiv0 ,
	\]
	即该向量场并不保守，不能直接断言积分为零。由 Green 公式，
	\[
	I=\iint\limits_{D}e^{xy}(y-x)\cos(x+y)\,dxdy .
	\]
	作变换 $(x,y)\mapsto(-x,-y)$：$D$ 关于原点对称保持不变，而
	$e^{xy}$ 不变、$\cos(x+y)$ 不变（余弦为偶函数）、$(y-x)$ 变号，
	故被积函数关于原点是奇函数，积分为零：
	\[
	I=0 .
	\]
}


\newpage
\item 
\begin{enumerate}
	%\renewcommand{\labelenumi}{\arabic{enumi}.}
	% A(\Alph) a(\alph) I(\Roman) i(\roman) 1(\arabic)
	%设定全局标号series=example	%引用全局变量resume=example
	%[topsep=-0.3em,parsep=-0.3em,itemsep=-0.3em,partopsep=-0.3em]
	%可使用leftmargin调整列表环境左边的空白长度 [leftmargin=0em]
	\item
	在 $[-\pi, \pi]$ 上把函数 $f(x)=x$ 展开为傅里叶级数；
	\item 
	证明: $\sum\limits_{n=0}^{+\infty} \frac{(-1)^{n}}{2 n+1}=\frac{\pi}{4}$。
	
\end{enumerate}


\banswer{
	\begin{enumerate}
		%\renewcommand{\labelenumi}{\arabic{enumi}.}
		% A(\Alph) a(\alph) I(\Roman) i(\roman) 1(\arabic)
		%设定全局标号series=example	%引用全局变量resume=example
		%[topsep=-0.3em,parsep=-0.3em,itemsep=-0.3em,partopsep=-0.3em]
		%可使用leftmargin调整列表环境左边的空白长度 [leftmargin=0em]
		\item
		$f(x)=x$ 在 $[-\pi,\pi]$ 上是奇函数，故 $a_n=0\ (n\geqslant0)$，只有正弦项。
		由
		\[
		b_n=\frac{1}{\pi}\int_{-\pi}^{\pi}x\sin nx\,dx=\frac{2}{\pi}\int_0^{\pi}x\sin nx\,dx ,
		\]
		分部积分，
		\[
		\int_0^{\pi}x\sin nx\,dx=\left[-\frac{x\cos nx}{n}\right]_0^{\pi}+\frac1n\int_0^{\pi}\cos nx\,dx
		=-\frac{\pi\cos n\pi}{n}+0=\frac{\pi(-1)^{n+1}}{n} ,
		\]
		故 $b_n=\frac{2(-1)^{n+1}}{n}$，从而
		\[
		f(x)=\sum\limits_{n=1}^{+\infty} \frac{2(-1)^{n+1}}{n} \sin n x ,\qquad x\in(-\pi,\pi).
		\]
		（在端点 $x=\pm\pi$ 处级数收敛到 $\frac{f(\pi^-)+f(-\pi^+)}{2}=0$，与 $f$ 的值不同。）
		\item 
		由 (1) 知 $x=\sum\limits_{n=1}^{+\infty}\frac{2(-1)^{n+1}}{n}\sin nx$，$x\in(-\pi,\pi)$。
		令 $x=\frac{\pi}{2}$，则 $\frac{\pi}{2}=\sum\limits_{n=1}^{+\infty}\frac{2(-1)^{n+1}}{n}\sin\frac{n\pi}{2}$。
		注意到 $\sin\frac{n\pi}{2}$ 在 $n$ 为偶数时为 $0$，$n=2k+1$ 时 $\sin\frac{(2k+1)\pi}{2}=(-1)^{k}$。
		故 $\frac{\pi}{2}=\sum\limits_{k=0}^{+\infty}\frac{2(-1)^{2k+2}}{2k+1}(-1)^{k}=2\sum\limits_{k=0}^{+\infty}\frac{(-1)^{k}}{2k+1}$，
		即 $\sum\limits_{n=0}^{+\infty}\frac{(-1)^{n}}{2n+1}=\frac{\pi}{4}$。
		
	\end{enumerate}
	
	
}





\item 
若$f(x)$和$ g(x)$ 在 $[a, b]$ 上满足二阶可导, 并且有 $g^{\prime \prime}(x) \neq 0$，$ f(a)=f(b)=g(a)=g(b)=0$，试证明：
\begin{enumerate}
	%\renewcommand{\labelenumi}{\arabic{enumi}.}
	% A(\Alph) a(\alph) I(\Roman) i(\roman) 1(\arabic)
	%设定全局标号series=example	%引用全局变量resume=example
	%[topsep=-0.3em,parsep=-0.3em,itemsep=-0.3em,partopsep=-0.3em]
	%可使用leftmargin调整列表环境左边的空白长度 [leftmargin=0em]
	\item
	$ \forall x \in(a, b)$，$  g(x) \neq 0 $；
	\item 
	$ \exists \xi \in(a, b)$，$ \frac{f(\xi)}{g(\xi)}=\frac{f^{\prime \prime}(\xi)}{g^{\prime \prime}(\xi)}$。
	
	
\end{enumerate}

\banswer{
	\begin{enumerate}
		\item
		假设存在 $x_{0}\in(a,b)$ 使 $g(x_{0})=0$。由 $g(a)=g(x_{0})=g(b)=0$，
		对 $g$ 在 $[a,x_{0}]$ 和 $[x_{0},b]$ 上分别应用 Rolle 定理，
		存在 $\eta_{1}\in(a,x_{0})$，$\eta_{2}\in(x_{0},b)$ 使 $g'(\eta_{1})=g'(\eta_{2})=0$。
		再对 $g'$ 在 $[\eta_{1},\eta_{2}]$ 上用 Rolle 定理，
		存在 $\xi\in(\eta_{1},\eta_{2})\subset(a,b)$ 使 $g''(\xi)=0$，与 $g''(x)\neq0$ 矛盾。
		故 $\forall x\in(a,b),\,g(x)\neq0$。\\
		\item
		令 $\varphi(x)=f(x)g'(x)-f'(x)g(x)$。由 $f(a)=g(a)=0$，$f(b)=g(b)=0$ 知
		$\varphi(a)=\varphi(b)=0$。由 Rolle 定理，$\exists\,\xi\in(a,b)$ 使 $\varphi'(\xi)=0$。
		计算 $\varphi'(x)=f'(x)g'(x)+f(x)g''(x)-f''(x)g(x)-f'(x)g'(x)=f(x)g''(x)-f''(x)g(x)$。
		由 $\varphi'(\xi)=0$ 得 $f(\xi)g''(\xi)=f''(\xi)g(\xi)$。
		由 (1) 知 $g(\xi)\neq0$，又 $g''(\xi)\neq0$，故
		$\displaystyle\frac{f(\xi)}{g(\xi)}=\frac{f''(\xi)}{g''(\xi)}$。
	\end{enumerate}
}




\item 
设$I_{n}=\int_{0}^{1} \frac{x^{n}}{1+x} d x$，证明：
\begin{enumerate}
	%\renewcommand{\labelenumi}{\arabic{enumi}.}
	% A(\Alph) a(\alph) I(\Roman) i(\roman) 1(\arabic)
	%设定全局标号series=example	%引用全局变量resume=example
	%[topsep=-0.3em,parsep=-0.3em,itemsep=-0.3em,partopsep=-0.3em]
	%可使用leftmargin调整列表环境左边的空白长度 [leftmargin=0em]
	\item %\label{2022-10-01}
	$I_{n+1}=-I_{n}+\frac{1}{n+1}$；
	\item 
	$\lim\limits _{n \rightarrow+\infty} I_{n}=0$；
	\item 
	用 $(1)$  和 $(2)$ 的结论证明 $\sum\limits_{n=1}^{+\infty} \frac{(-1)^{n-1}}{n}=\ln 2$。
	
\end{enumerate}

\banswer{
	\begin{enumerate}
		\item
		$I_{n}+I_{n+1}=\int_{0}^{1}\frac{x^{n}+x^{n+1}}{1+x}\,dx
		=\int_{0}^{1}\frac{x^{n}(1+x)}{1+x}\,dx
		=\int_{0}^{1}x^{n}\,dx=\frac{1}{n+1}$，
		故 $I_{n+1}=-I_{n}+\frac{1}{n+1}$。\\
		\item
		$0\leq I_{n}=\int_{0}^{1}\frac{x^{n}}{1+x}\,dx
		\leq\int_{0}^{1}x^{n}\,dx=\frac{1}{n+1}\to0\;(n\to\infty)$，
		由夹逼定理 $\lim\limits_{n\to+\infty}I_{n}=0$。\\
		\item
		由 (1) 递推得 $I_{n}=(-1)^{n}\bigl(I_{0}-\sum_{k=1}^{n}\frac{(-1)^{k-1}}{k}\bigr)$。
		其中 $I_{0}=\int_{0}^{1}\frac{1}{1+x}\,dx=\ln2$。
		令 $n\to\infty$，由 (2) $I_{n}\to0$，得
		$\displaystyle\ln2-\sum_{k=1}^{\infty}\frac{(-1)^{k-1}}{k}=0$，
		即 $\displaystyle\sum_{n=1}^{+\infty}\frac{(-1)^{n-1}}{n}=\ln2$。
	\end{enumerate}
}





\item 
函数 $f(x, y, z)=x^{2}+y^{2}+z^{2}$, 求其在满足条件 $a x+b y+c z=1$ 下的最小值。
\banswer{
	设 $\bm{n}=(a,b,c)\neq\bm{0}$，$\bm{r}=(x,y,z)$，约束即 $\bm{n}\cdot\bm{r}=1$。

	\textbf{解法一（Cauchy--Schwarz）.} 由
	\[
	1=(ax+by+cz)^{2}\leqslant\left(a^{2}+b^{2}+c^{2}\right)\left(x^{2}+y^{2}+z^{2}\right)
	\]
	立得 $f=x^{2}+y^{2}+z^{2}\geqslant\frac{1}{a^{2}+b^{2}+c^{2}}$，
	等号当且仅当 $\bm{r}\parallel\bm{n}$ 时取到。

	\textbf{解法二（Lagrange 乘数法）.} 令
	$L=x^{2}+y^{2}+z^{2}-\lambda(ax+by+cz-1)$，由
	$2x=\lambda a$、$2y=\lambda b$、$2z=\lambda c$ 得 $\bm{r}=\frac{\lambda}{2}\bm{n}$；
	代入约束得 $\frac{\lambda}{2}\left(a^{2}+b^{2}+c^{2}\right)=1$，故
	\[
	\bm{r}=\frac{(a,b,c)}{a^{2}+b^{2}+c^{2}},\qquad
	f_{\min}=|\bm{r}|^{2}=\frac{1}{a^{2}+b^{2}+c^{2}} .
	\]
	（几何意义：$f$ 的最小值就是原点到平面 $ax+by+cz=1$ 的距离的平方
	$\left(\frac{1}{\sqrt{a^{2}+b^{2}+c^{2}}}\right)^{2}$。）
}

	
\end{enumerate}

